\section{Effectful Mealy morphisms}
\label{sec:effectful-mealy-morphisms}

The Mealy morphisms considered so far are deterministic. For every admissible
input-state pair
\[
    (a,x)\in A_1\times_{t_A,A_0,p}P,
\]
the structure map selects a single updated state
\[
    \delta(a,x)
\]
and a single output arrow
\[
    \omega(a,x):
    q(\delta(a,x))\to q(x).
\]
Equivalently, it selects a single element of the apex of
\[
    T_B(P)=B\circ P.
\]

To introduce computational effects, the next state and the output arrow should
be placed jointly inside an effect. Thus an effectful transition has the form
\[
    P\circ A
    \longrightarrow
    \mathsf K(B\circ P),
\]
where \(\mathsf K\) is a monad preserving the external interfaces of the
carrier. The order of the composite is important: the output arrow is first
attached to the updated state, and the effect is then applied to the resulting
joint state-output object. Consequently, different effectful branches may have
different next states and different output arrows.

The intended execution shape is
\[
    \begin{tikzcd}[column sep=huge]
            u
                \arrow[r, "a"]
        &
            v=p(x)
    \end{tikzcd}
    \qquad\leadsto\qquad
    \mathsf K
    \left[
        \begin{tikzcd}[column sep=huge, baseline=-0.5ex]
                q(x')
                    \arrow[r, "\beta"]
            &
                q(x)
        \end{tikzcd}
    \right].
\]
Here both \(x'\) and \(\beta\) occur inside the same computational effect.

Related categorical approaches to effectful Mealy machines have been developed
by Bonchi, Di Lavore, and Román
\cite{BonchiDiLavoreRomanEffectfulMealyMachines}. The construction below is
adapted to the present setting of internal categories and span-carried state.

\subsection{Interface-preserving effects}
\label{subsec:interface-preserving-effects}

For objects \(X,Y\in\mathcal E\), write
\[
    \mathcal C_{X,Y}
    :=
    \mathbf{Span}(\mathcal E)(X,Y).
\]
Since \(\mathcal E\) has finite products, there is a canonical isomorphism of
categories
\[
    \mathcal C_{X,Y}
    \cong
    \mathcal E/(X\times Y).
\]
It sends a span
\[
    X\xleftarrow{p}P\xrightarrow{q}Y
\]
to the morphism
\[
    \langle p,q\rangle:P\to X\times Y.
\]
The correspondence is displayed by
\[
    \begin{tikzcd}[column sep=huge]
            &
                P
                    \arrow[dl, "p"']
                    \arrow[dr, "q"]
                    \arrow[d, "{\langle p,q\rangle}" description]
            &
        \\
            X
            &
                X\times Y
                    \arrow[l, "\pi_X"]
                    \arrow[r, "\pi_Y"']
            &
                Y.
    \end{tikzcd}
\]

A monad on \(\mathcal C_{X,Y}\), equivalently on
\(\mathcal E/(X\times Y)\), preserves the external interfaces in the sense
that every effectful carrier remains equipped with a specified map to
\(X\times Y\). In the fibrewise examples below, a stronger property holds:
the effect acts pointwise and independently on each interface fibre. The
abstract definition does not require this stronger pointwise condition.

Fix internal categories
\[
    \mathbb A=(A_0,A_1,s_A,t_A,e_A,m_A),
    \qquad
    \mathbb B=(B_0,B_1,s_B,t_B,e_B,m_B).
\]
Recall the local monads
\[
    R_A=-\circ A,
    \qquad
    T_B=B\circ-
\]
on
\[
    \mathcal C_{A,B}
    :=
    \mathbf{Span}(\mathcal E)(A_0,B_0),
\]
together with the canonical distributive law
\[
    \lambda:
    R_AT_B
    \Rightarrow
    T_BR_A
\]
from Proposition~\ref{prop:canonical-distributive-law}.

\begin{definition}[Mealy-compatible effect datum]
\label{def:mealy-compatible-effect-datum}
    A \textbf{Mealy-compatible effect datum} for
    \[
        (\mathbb A,\mathbb B)
    \]
    consists of the following structure on
    \(\mathcal C_{A,B}\).

    \begin{enumerate}
        \item A monad
        \[
            \mathsf K
            =
            (\mathsf K,\eta^{\mathsf K},\mu^{\mathsf K}).
        \]

        \item A distributive law
        \[
            \sigma_B:
            T_B\mathsf K
            \Rightarrow
            \mathsf K T_B
        \]
        of the output-comparison monad over the effect monad.

        \item A distributive law
        \[
            \rho_A:
            R_A\mathsf K
            \Rightarrow
            \mathsf K R_A
        \]
        of the input-processing monad over the effect monad.

        \item The Yang--Baxter compatibility equation
        \begin{equation}
        \begin{aligned}
            &
            (\mathsf K\lambda)
            \circ
            (\rho_A T_B)
            \circ
            (R_A\sigma_B)
            \\
            &\qquad=
            (\sigma_B R_A)
            \circ
            (T_B\rho_A)
            \circ
            (\lambda\mathsf K)
        \end{aligned}
        \label{eq:effect-yang-baxter}
        \end{equation}
        as natural transformations
        \[
            R_AT_B\mathsf K
            \Rightarrow
            \mathsf K T_BR_A.
        \]
    \end{enumerate}
\end{definition}

The Yang--Baxter equation is the commutativity of the diagram
\[
    \begin{tikzcd}[column sep=huge, row sep=large]
            R_AT_B\mathsf K
                \arrow[r, "R_A\sigma_B"]
                \arrow[d, "\lambda\mathsf K"']
        &
            R_A\mathsf K T_B
                \arrow[r, "\rho_A T_B"]
        &
            \mathsf K R_AT_B
                \arrow[d, "\mathsf K\lambda"]
        \\
            T_BR_A\mathsf K
                \arrow[r, "T_B\rho_A"']
        &
            T_B\mathsf K R_A
                \arrow[r, "\sigma_B R_A"']
        &
            \mathsf K T_BR_A.
    \end{tikzcd}
\]
It says that input processing, output comparison, and the computational effect
may be interchanged in either order with the same result.

\subsection{The canonical fibrewise effect doctrine in
\texorpdfstring{\(\mathbf{Set}\)}{Set}}
\label{subsec:fibrewise-effects-set}

Let
\[
    \mathcal E=\mathbf{Set},
\]
and let
\[
    K=(K,\eta^K,\mu^K)
\]
be a monad on \(\mathbf{Set}\).

For a span
\[
    X\xleftarrow{p}P\xrightarrow{q}Y,
\]
write
\[
    P_{x,y}
    :=
    \{z\in P\mid p(z)=x,\ q(z)=y\}.
\]
Thus
\[
    P
    \cong
    \coprod_{(x,y)\in X\times Y}P_{x,y}.
\]

Define
\[
    \mathsf K_{X,Y}(P)
    :=
    \coprod_{(x,y)\in X\times Y}K(P_{x,y}),
\]
with the evident projection
\[
    \mathsf K_{X,Y}(P)\to X\times Y.
\]
The resulting carrier span is
\[
    \begin{tikzcd}
        &
            \displaystyle
            \coprod_{(x,y)\in X\times Y}K(P_{x,y})
                \arrow[dl]
                \arrow[dr]
        &
        \\
            X
        &
        &
            Y.
    \end{tikzcd}
\]
If
\[
    f:P\to Q
\]
is a map of spans, then \(f\) restricts to maps
\[
    f_{x,y}:P_{x,y}\to Q_{x,y},
\]
and
\[
    \mathsf K_{X,Y}(f)
    =
    \coprod_{(x,y)}
    K(f_{x,y}).
\]

\begin{proposition}[Fibrewise effect doctrine in
\texorpdfstring{\(\mathbf{Set}\)}{Set}]
\label{prop:fibrewise-effect-doctrine-set}
    Every monad
    \[
        K=(K,\eta^K,\mu^K)
    \]
    on \(\mathbf{Set}\) induces, for each pair of sets \(X,Y\), a monad
    \[
        \mathsf K_{X,Y}
    \]
    on
    \[
        \mathbf{Span}(\mathbf{Set})(X,Y)
        \cong
        \mathbf{Set}/(X\times Y)
    \]
    by applying \(K\) separately to each interface fibre.

    For fixed small categories \(\mathbb A\) and \(\mathbb B\), this monad
    carries canonical distributive laws
    \[
        \sigma_B:
        T_B\mathsf K_{A_0,B_0}
        \Rightarrow
        \mathsf K_{A_0,B_0}T_B
    \]
    and
    \[
        \rho_A:
        R_A\mathsf K_{A_0,B_0}
        \Rightarrow
        \mathsf K_{A_0,B_0}R_A
    \]
    satisfying the Yang--Baxter equation
    \eqref{eq:effect-yang-baxter}.
\end{proposition}
\begin{proof}
    \leavevmode
    \begin{enumerate}
        \item \textbf{Construct the fibrewise monad.}
        The unit and multiplication are defined fibrewise:
        \[
            \eta^{\mathsf K}_P
            =
            \coprod_{(x,y)}
            \eta^K_{P_{x,y}},
        \]
        \[
            \mu^{\mathsf K}_P
            =
            \coprod_{(x,y)}
            \mu^K_{P_{x,y}}.
        \]
        The monad equations hold because they hold separately on every
        interface fibre.

        \item \textbf{Construct the output-comparison law.}
        Let
        \[
            P=
            \left(
                A_0\xleftarrow{p}P\xrightarrow{q}B_0
            \right).
        \]
        A generalized element of \(T_B\mathsf K(P)\) consists of an
        effectful state
        \[
            \kappa\in K(P_{u,c})
        \]
        and an arrow
        \[
            \beta:c\to b
        \]
        in \(\mathbb B\):
        \[
            \begin{tikzcd}[column sep=huge]
                    c
                        \arrow[r, "\beta"]
                &
                    b.
            \end{tikzcd}
        \]
        The map
        \[
            P_{u,c}
            \longrightarrow
            (T_BP)_{u,b},
            \qquad
            x\longmapsto(x,\beta),
        \]
        induces
        \[
            K(P_{u,c})
            \longrightarrow
            K((T_BP)_{u,b}).
        \]
        Define
        \[
            \sigma_{B,P}(\kappa,\beta)
            :=
            K\bigl(x\mapsto(x,\beta)\bigr)(\kappa).
        \]
        Diagrammatically, the output arrow is moved inside the effect:
        \[
            \begin{tikzcd}[column sep=huge]
                    \kappa\in K(P_{u,c})
                        \arrow[r, maps to]
                &
                    K\bigl\{
                        (x,\beta)
                        \mid
                        x\in P_{u,c}
                    \bigr\}.
            \end{tikzcd}
        \]

        \item \textbf{Construct the input-processing law.}
        A generalized element of \(R_A\mathsf K(P)\) consists of an arrow
        \[
            a:u\to v
        \]
        and an effectful state
        \[
            \kappa\in K(P_{v,b}).
        \]
        The map
        \[
            P_{v,b}
            \longrightarrow
            (R_AP)_{u,b},
            \qquad
            x\longmapsto(a,x),
        \]
        induces
        \[
            K(P_{v,b})
            \longrightarrow
            K((R_AP)_{u,b}).
        \]
        Define
        \[
            \rho_{A,P}(a,\kappa)
            :=
            K\bigl(x\mapsto(a,x)\bigr)(\kappa).
        \]
        The input arrow is therefore moved inside the effect:
        \[
            \begin{tikzcd}[column sep=huge]
                    u
                        \arrow[r, "a"]
                &
                    v
                        \arrow[r, phantom, "\kappa\in K(P_{v,b})"]
                &
                    {}
            \end{tikzcd}
            \qquad\leadsto\qquad
            K\bigl\{(a,x)\mid x\in P_{v,b}\bigr\}.
        \]

        \item \textbf{Verify the distributive-law equations.}
        The unit equations follow from naturality of \(\eta^K\). The
        multiplication equations follow from naturality of \(\mu^K\).
        For example, attaching an output arrow before or after flattening a
        nested effect gives the same map because
        \[
            \mu^K
            \circ
            KK(x\mapsto(x,\beta))
            =
            K(x\mapsto(x,\beta))
            \circ
            \mu^K.
        \]
        The input equations are identical in form.

        \item \textbf{Verify Yang--Baxter compatibility.}
        Fix
        \[
            a:u\to v,
            \qquad
            \beta:c\to b,
            \qquad
            \kappa\in K(P_{v,c}).
        \]
        Both sides of the Yang--Baxter equation apply \(K\) to the same map
        \[
            P_{v,c}
            \longrightarrow
            (T_BR_AP)_{u,b},
            \qquad
            x\longmapsto((a,x),\beta).
        \]
        The two presentations differ only by the canonical reassociation
        \[
            (a,(x,\beta))
            \longleftrightarrow
            ((a,x),\beta).
        \]
        Hence the diagram
        \[
            \begin{tikzcd}[column sep=huge, row sep=large]
                    (a,\kappa,\beta)
                        \arrow[r, maps to]
                        \arrow[d, maps to]
                &
                    (a,K(x\mapsto(x,\beta))(\kappa))
                        \arrow[d, maps to]
                \\
                    ((a,\kappa),\beta)
                        \arrow[r, maps to]
                &
                    K(x\mapsto((a,x),\beta))(\kappa)
            \end{tikzcd}
        \]
        commutes. This is precisely
        \eqref{eq:effect-yang-baxter}.
    \end{enumerate}
\end{proof}

Important examples of the fibrewise construction are the finite-powerset
monad, the full powerset monad, the finite-distribution monad, the
finite-subdistribution monad, the lift monad, and finitely supported weighted
multiset monads.

\subsection{The effectful output-comparison monad}
\label{subsec:effectful-output-comparison-monad}

The output-comparison monad and the computational effect combine into a single
monad.

\begin{proposition}[The effectful output-comparison monad]
\label{prop:effectful-output-comparison-monad}
    Let
    \[
        (\mathsf K,\sigma_B,\rho_A)
    \]
    be a Mealy-compatible effect datum. The composite endofunctor
    \[
        H_B^{\mathsf K}
        :=
        \mathsf K T_B
    \]
    carries a canonical monad structure
    \[
        H_B^{\mathsf K}
        =
        (\mathsf K T_B,\eta^H,\mu^H)
    \]
    on \(\mathcal C_{A,B}\).
\end{proposition}
\begin{proof}
    \leavevmode
    \begin{enumerate}
        \item \textbf{Use the output-effect distributive law.}
        The transformation
        \[
            \sigma_B:
            T_B\mathsf K
            \Rightarrow
            \mathsf K T_B
        \]
        is a distributive law of \(T_B\) over \(\mathsf K\). Therefore the
        composite endofunctor
        \[
            \mathsf K T_B
        \]
        carries the standard composite-monad structure.

        \item \textbf{Record the unit.}
        The unit is
        \[
            \eta^H_P:
            P
            \xrightarrow{\eta^T_P}
            T_B(P)
            \xrightarrow{\eta^{\mathsf K}_{T_B(P)}}
            \mathsf K T_B(P).
        \]
        Thus
        \[
            x
            \longmapsto
            \eta^{\mathsf K}
            \bigl(x,e_B(qx)\bigr).
        \]
        The corresponding output-comparison diagram is
        \[
            \begin{tikzcd}[column sep=huge]
                    q(x)
                        \arrow[r, "e_B(qx)"]
                &
                    q(x).
            \end{tikzcd}
        \]

        \item \textbf{Record the multiplication.}
        The multiplication is
        \[
        \begin{aligned}
            \mu^H_P:
            \mathsf K T_B\mathsf K T_B(P)
                &\xrightarrow{\mathsf K\sigma_B T_B}
            \mathsf K\mathsf K T_BT_B(P)
            \\
                &\xrightarrow{\mu^{\mathsf K}T_BT_B}
            \mathsf K T_BT_B(P)
            \xrightarrow{\mathsf K\mu^T_P}
            \mathsf K T_B(P).
        \end{aligned}
        \]
        The effect multiplication combines the nested computational effects,
        while \(\mu^T\) composes the output-comparison arrows:
        \[
            \begin{tikzcd}[column sep=huge]
                    q(x)
                        \arrow[r, "\beta"]
                &
                    b
                        \arrow[r, "\gamma"]
                &
                    c
            \end{tikzcd}
            \qquad\longmapsto\qquad
            \begin{tikzcd}[column sep=huge]
                    q(x)
                        \arrow[r, "\gamma\circ\beta"]
                &
                    c.
            \end{tikzcd}
        \]

        \item \textbf{Apply the composite-monad theorem.}
        The monad laws follow from the monad laws for \(T_B\) and
        \(\mathsf K\), together with the distributive-law equations for
        \(\sigma_B\).
    \end{enumerate}
\end{proof}

A value of
\[
    H_B^{\mathsf K}(P)
    =
    \mathsf K T_B(P)
\]
is an effectful value whose pure outcomes are pairs
\[
    (x,\beta)
\]
with
\[
    \beta:q(x)\to b.
\]
The state and output-comparison arrow occur inside the same effect and may
therefore be correlated.

\begin{remark}[Order of the composite]
\label{rem:effectful-output-composite-order}
    The intended monad is
    \[
        \mathsf K T_B,
    \]
    not
    \[
        T_B\mathsf K.
    \]
    A pure outcome of \(\mathsf K T_B(P)\) is a joint pair
    \[
        (x,\beta).
    \]
    Thus distinct nondeterministic or probabilistic branches may produce
    distinct next states and distinct output arrows.

    By contrast, an element of \(T_B\mathsf K(P)\) consists of an effectful
    state together with one output arrow outside the effect:
    \[
        \begin{tikzcd}[column sep=huge]
                \mathsf K\{x\}
                    \arrow[r, phantom, "\text{one fixed }\beta"]
            &
                {}
        \end{tikzcd}
    \]
    This would force every branch to share the same output-comparison arrow.
\end{remark}

\subsection{Lifting input processing through the effect}
\label{subsec:effectful-input-processing-lift}

The three pairwise distributive laws combine by the iterated
distributive-law theorem.

\begin{proposition}[The effectful input distributive law]
\label{prop:effectful-input-distributive-law}
    Define
    \[
        \chi:
        R_AH_B^{\mathsf K}
        \Rightarrow
        H_B^{\mathsf K}R_A
    \]
    by
    \[
        \chi
        :=
        (\mathsf K\lambda)
        \circ
        (\rho_A T_B).
    \]
    Thus the component at \(P\) is
    \[
        R_A\mathsf K T_B(P)
        \xrightarrow{\rho_A T_B}
        \mathsf K R_AT_B(P)
        \xrightarrow{\mathsf K\lambda}
        \mathsf K T_BR_A(P).
    \]
    Then \(\chi\) is a distributive law of the monad \(R_A\) over the monad
    \(H_B^{\mathsf K}\).
\end{proposition}
\begin{proof}
    The monads
    \[
        R_A,
        \qquad
        T_B,
        \qquad
        \mathsf K
    \]
    are equipped with distributive laws
    \[
        \lambda:
        R_AT_B\Rightarrow T_BR_A,
    \]
    \[
        \rho_A:
        R_A\mathsf K\Rightarrow\mathsf K R_A,
    \]
    and
    \[
        \sigma_B:
        T_B\mathsf K\Rightarrow\mathsf K T_B.
    \]
    By Definition~\ref{def:mealy-compatible-effect-datum}, these laws satisfy
    the Yang--Baxter equation
    \eqref{eq:effect-yang-baxter}.

    The iterated distributive-law theorem therefore applies
    \cite{ChengIteratedDistributiveLaws}. It implies simultaneously that
    \[
        H_B^{\mathsf K}
        =
        \mathsf K T_B
    \]
    is a monad and that
    \[
        R_A\mathsf K T_B
        \xrightarrow{\rho_A T_B}
        \mathsf K R_AT_B
        \xrightarrow{\mathsf K\lambda}
        \mathsf K T_BR_A
    \]
    is a distributive law of \(R_A\) over \(\mathsf K T_B\).

    The transformation is displayed by
    \[
        \begin{tikzcd}[column sep=huge]
                R_A\mathsf K T_B(P)
                    \arrow[r, "\rho_A T_B"]
            &
                \mathsf K R_AT_B(P)
                    \arrow[r, "\mathsf K\lambda"]
            &
                \mathsf K T_BR_A(P).
        \end{tikzcd}
    \]
    The distributive-law equations are exactly the pairwise distributive-law
    equations together with the Yang--Baxter compatibility.
\end{proof}

\begin{corollary}[The effectful lifted input monad]
\label{cor:effectful-lifted-input-monad}
    The distributive law of
    Proposition~\ref{prop:effectful-input-distributive-law} lifts \(R_A\) to a
    monad
    \[
        \overline R_A^{\mathsf K}
    \]
    on the Kleisli category
    \[
        \operatorname{Kl}(H_B^{\mathsf K}).
    \]
\end{corollary}
\begin{proof}
    This is the Kleisli lifting theorem applied to
    \[
        \chi:
        R_AH_B^{\mathsf K}
        \Rightarrow
        H_B^{\mathsf K}R_A.
    \]
    On objects,
    \[
        \overline R_A^{\mathsf K}(P)=R_A(P).
    \]
    If a Kleisli arrow
    \[
        f:P\to Q
    \]
    is represented in \(\mathcal C_{A,B}\) by
    \[
        f:P\to H_B^{\mathsf K}(Q),
    \]
    then its image is represented by
    \[
        R_A(P)
        \xrightarrow{R_Af}
        R_AH_B^{\mathsf K}(Q)
        \xrightarrow{\chi_Q}
        H_B^{\mathsf K}R_A(Q).
    \]
\end{proof}

\subsection{Effectful Mealy morphisms}
\label{subsec:effectful-mealy-morphisms-definition}

\begin{definition}[Effectful Mealy morphism]
\label{def:effectful-mealy-morphism}
    Let
    \[
        (\mathsf K,\sigma_B,\rho_A)
    \]
    be a Mealy-compatible effect datum for
    \[
        (\mathbb A,\mathbb B).
    \]

    A \textbf{\(\mathsf K\)-effectful Mealy morphism}
    \[
        (P,\widehat\Phi):
        \mathbb A\rightsquigarrow_{\mathsf K}\mathbb B
    \]
    consists of a carrier span
    \[
        A_0\xleftarrow{p}P\xrightarrow{q}B_0
    \]
    and an algebra structure for
    \[
        \overline R_A^{\mathsf K}
    \]
    on \(P\) in
    \[
        \operatorname{Kl}(H_B^{\mathsf K}).
    \]

    Equivalently, it consists of a map in
    \(\mathcal C_{A,B}\)
    \[
        \widehat\Phi:
        R_A(P)
        \longrightarrow
        H_B^{\mathsf K}(P),
    \]
    that is,
    \[
        \widehat\Phi:
        P\circ A
        \longrightarrow
        \mathsf K(B\circ P),
    \]
    satisfying
    \begin{align}
        \widehat\Phi\circ\eta^R_P
        &=
        \eta^H_P,
        \label{eq:effectful-mealy-unit-abstract}
        \\
        \widehat\Phi\circ\mu^R_P
        &=
        \mu^H_P
        \circ
        H_B^{\mathsf K}(\widehat\Phi)
        \circ
        \chi_P
        \circ
        R_A(\widehat\Phi).
        \label{eq:effectful-mealy-multiplication-abstract}
    \end{align}
\end{definition}

The structure map is a morphism of carrier spans
\[
    \begin{tikzcd}[column sep=huge]
        &
            A_1\times_{t_A,A_0,p}P
                \arrow[d, "\widehat\Phi" description]
                \arrow[dl, "s_A\pi_{A_1}"']
                \arrow[dr, "q\pi_P"]
        &
        \\
            A_0
        &
            \mathsf K(B\circ P)
                \arrow[l]
                \arrow[r]
        &
            B_0.
    \end{tikzcd}
\]
Thus the effect does not change the external input and output interfaces of the
transition.

\subsection{Fibrewise component laws}
\label{subsec:fibrewise-effectful-mealy-laws}

Assume throughout this subsection that
\[
    \mathcal E=\mathbf{Set}
\]
and that the effect datum is the canonical fibrewise datum induced by a
Set-monad
\[
    K=(K,\eta^K,\mu^K)
\]
as in Proposition~\ref{prop:fibrewise-effect-doctrine-set}.

For \(u\in A_0\) and \(b\in B_0\), define the set of typed output outcomes
\[
    \operatorname{Out}_P(u,b)
    :=
    \left\{
        (x',\beta)\in P\times B_1
        \ \middle|\
        \begin{aligned}
            p(x')&=u,\\
            s_B(\beta)&=q(x'),\\
            t_B(\beta)&=b
        \end{aligned}
    \right\}.
\]
This is the fibre of \(T_B(P)\) over \((u,b)\).

An admissible input-state pair consists of
\[
    a:u\to v,
    \qquad
    x\in P,
    \qquad
    p(x)=v.
\]
The effectful transition has type
\[
    \widehat\Phi(a,x)
    \in
    K\bigl(
        \operatorname{Out}_P(u,q(x))
    \bigr).
\]
Every pure outcome
\[
    (x',\beta)
\]
has the typing diagram
\[
    \begin{tikzcd}[column sep=huge, row sep=large]
            u
                \arrow[r, "a"]
        &
            v=p(x)
        \\
            q(x')
                \arrow[r, "\beta"']
        &
            q(x),
    \end{tikzcd}
\]
where
\[
    p(x')=u.
\]

For a Set-monad \(K\), write
\[
    \operatorname{bind}_{K}(\kappa,f)
    :=
    \mu^K\bigl(K(f)(\kappa)\bigr).
\]

\begin{proposition}[The fibrewise effectful Mealy laws]
\label{prop:effectful-mealy-laws}
    Under the canonical fibrewise Set effect datum, the abstract equations
    \eqref{eq:effectful-mealy-unit-abstract} and
    \eqref{eq:effectful-mealy-multiplication-abstract} are equivalent to the
    following component equations.

    The unit law is
    \begin{equation}
        \widehat\Phi(e_A(px),x)
        =
        \eta^K
        \bigl(x,e_B(qx)\bigr).
        \label{eq:effectful-mealy-unit-component}
    \end{equation}

    For internally composable arrows
    \[
        a:u\to v,
        \qquad
        b:v\to w,
    \]
    and a state \(x\in P\) satisfying \(p(x)=w\), the multiplication law is
    \begin{equation}
    \begin{aligned}
        \widehat\Phi(b\circ a,x)
        =
        \operatorname{bind}_{K}
        \Bigl(
            \widehat\Phi(b,x),
            (x_b,\beta_b)
            \longmapsto
            \operatorname{bind}_{K}
            \bigl(
                \widehat\Phi(a,x_b),
                (x_a,\beta_a)
                \longmapsto
                \eta^K
                (x_a,\beta_b\circ\beta_a)
            \bigr)
        \Bigr).
    \end{aligned}
    \label{eq:effectful-mealy-composition-component}
    \end{equation}
\end{proposition}
\begin{proof}
    \leavevmode
    \begin{enumerate}
        \item \textbf{Unpack the unit.}
        The unit of \(R_A\) sends
        \[
            x
            \longmapsto
            (e_A(px),x).
        \]
        The unit of \(H_B^{\mathsf K}\) attaches the identity output arrow and
        inserts the result as a pure \(K\)-outcome:
        \[
            x
            \longmapsto
            \eta^K
            \bigl(x,e_B(qx)\bigr).
        \]
        Hence the algebra unit equation is precisely
        \eqref{eq:effectful-mealy-unit-component}.

        \item \textbf{Unpack the first execution.}
        The sequential side first executes
        \[
            b:v\to w
        \]
        at \(x\). It produces effectfully a pair
        \[
            (x_b,\beta_b)
        \]
        with
        \[
            p(x_b)=v,
        \]
        and
        \[
            \beta_b:
            q(x_b)\to q(x).
        \]
        The output is displayed by
        \[
            \begin{tikzcd}[column sep=huge]
                    q(x_b)
                        \arrow[r, "\beta_b"]
                &
                    q(x).
            \end{tikzcd}
        \]

        \item \textbf{Unpack the second execution.}
        Since \(p(x_b)=v=t_A(a)\), the arrow
        \[
            a:u\to v
        \]
        is admissible at \(x_b\). Its execution produces effectfully
        \[
            (x_a,\beta_a)
        \]
        with
        \[
            p(x_a)=u,
        \]
        and
        \[
            \beta_a:
            q(x_a)\to q(x_b).
        \]

        \item \textbf{Compose the emitted arrows.}
        The output-comparison multiplication composes
        \[
            \begin{tikzcd}[column sep=huge]
                    q(x_a)
                        \arrow[r, "\beta_a"]
                &
                    q(x_b)
                        \arrow[r, "\beta_b"]
                &
                    q(x)
            \end{tikzcd}
        \]
        to obtain
        \[
            \beta_b\circ\beta_a:
            q(x_a)\to q(x).
        \]
        The multiplication of \(K\) combines the two nested effects by
        Kleisli extension.

        \item \textbf{Compare direct and sequential execution.}
        The other side first composes the input arrows:
        \[
            \begin{tikzcd}[column sep=huge]
                    u
                        \arrow[r, "a"]
                &
                    v
                        \arrow[r, "b"]
                &
                    w=p(x).
            \end{tikzcd}
        \]
        It then executes
        \[
            b\circ a:u\to w
        \]
        directly at \(x\). Equality of the resulting \(K\)-values is exactly
        \eqref{eq:effectful-mealy-composition-component}.
    \end{enumerate}
\end{proof}

\begin{remark}[Correlated state and output]
\label{rem:effectful-correlated-state-output}
    An effectful Mealy morphism does not separately select an effectful next
    state and an effectful output. It selects an effectful joint pair
    \[
        (x',\beta).
    \]
    Consequently, the next state and emitted output may be correlated. In the
    probabilistic case, observing the output may therefore reveal information
    about which state transition occurred.
\end{remark}

\subsection{Pure effectful Mealy morphisms}
\label{subsec:pure-effectful-mealy-morphisms}

The deterministic theory embeds into the effectful theory at the abstract
level.

Define
\[
    j:
    T_B
    \Rightarrow
    H_B^{\mathsf K}
\]
by
\[
    j
    :=
    \eta^{\mathsf K}T_B.
\]
Thus
\[
    j_P:
    T_B(P)
    \longrightarrow
    \mathsf K T_B(P)
\]
inserts a deterministic state-output pair as a pure effectful outcome.

\begin{proposition}[Pure effectful Mealy morphisms]
\label{prop:pure-effectful-mealy-morphisms}
    Every ordinary Mealy morphism
    \[
        (P,\Phi):
        \mathbb A\rightsquigarrow\mathbb B
    \]
    determines a canonical \(\mathsf K\)-effectful Mealy morphism with the
    same carrier, defined by
    \[
        \widehat\Phi
        :=
        j_P\circ\Phi
        =
        \eta^{\mathsf K}_{T_B(P)}
        \circ
        \Phi.
    \]
\end{proposition}
\begin{proof}
    The distributive-law axioms for \(\sigma_B\) imply that
    \[
        j:
        T_B\Rightarrow H_B^{\mathsf K}
    \]
    is a morphism of monads. It therefore induces a functor
    \[
        J:
        \operatorname{Kl}(T_B)
        \longrightarrow
        \operatorname{Kl}(H_B^{\mathsf K})
    \]
    which is the identity on objects and sends
    \[
        f:X\to T_B(Y)
    \]
    to
    \[
        j_Yf:
        X\to H_B^{\mathsf K}(Y).
    \]

    This functor is compatible with the lifted input-processing monads. Indeed,
    the diagram
    \[
        \begin{tikzcd}[column sep=huge, row sep=large]
                R_AT_B
                    \arrow[r, "\lambda"]
                    \arrow[d, "R_Aj"']
            &
                T_BR_A
                    \arrow[d, "jR_A"]
            \\
                R_AH_B^{\mathsf K}
                    \arrow[r, "\chi"']
            &
                H_B^{\mathsf K}R_A
        \end{tikzcd}
    \]
    commutes. Algebraically,
    \[
    \begin{aligned}
        \chi\circ R_Aj
        &=
        (\mathsf K\lambda)
        \circ
        (\rho_AT_B)
        \circ
        R_A(\eta^{\mathsf K}T_B)
        \\
        &=
        (\mathsf K\lambda)
        \circ
        (\eta^{\mathsf K}R_AT_B)
        \\
        &=
        (\eta^{\mathsf K}T_BR_A)
        \circ
        \lambda
        \\
        &=
        (jR_A)\circ\lambda.
    \end{aligned}
    \]
    The second equality is the \(\mathsf K\)-unit equation for \(\rho_A\),
    and the third is naturality of \(\eta^{\mathsf K}\).

    Hence \(J\) sends \(\overline R_A\)-algebras to
    \(\overline R_A^{\mathsf K}\)-algebras. Applying \(J\) to
    \[
        \Phi:
        R_A(P)\to T_B(P)
    \]
    gives
    \[
        \widehat\Phi
        =
        \eta^{\mathsf K}_{T_B(P)}\Phi:
        R_A(P)\to H_B^{\mathsf K}(P),
    \]
    as required.
\end{proof}

In the fibrewise Set case, this sends
\[
    \Phi(a,x)
    =
    \bigl(\delta(a,x),\omega(a,x)\bigr)
\]
to
\[
    \widehat\Phi(a,x)
    =
    \eta^K
    \bigl(\delta(a,x),\omega(a,x)\bigr).
\]

\subsection{The one-object specialization}
\label{subsec:one-object-effectful-mealy-morphisms}

Let \(M\) and \(N\) be monoid objects, regarded as one-object internal
categories using the standing convention
\[
    b\circ a=ba.
\]
Since
\[
    \mathcal C_{1,1}
    =
    \mathbf{Span}(\mathcal E)(1,1)
    \cong
    \mathcal E,
\]
an abstract one-object effect datum consists in particular of a monad
\[
    \mathsf K:\mathcal E\to\mathcal E
\]
together with distributive laws governing its interaction with
\[
    M\times-
    \qquad\text{and}\qquad
    -\times N.
\]
An effectful one-object Mealy morphism has an underlying map
\[
    \widehat\Phi:
    M\times P
    \longrightarrow
    \mathsf K(P\times N).
\]

We now specialize to the canonical fibrewise datum induced by a Set-monad
\(K\).

\begin{proposition}[One-object fibrewise effectful Mealy morphisms]
\label{prop:one-object-effectful-mealy-morphisms}
    Let \(K\) be a monad on \(\mathbf{Set}\). A one-object
    \(K\)-effectful Mealy morphism
    \[
        M\rightsquigarrow_KN
    \]
    is equivalently a set \(P\) and a map
    \[
        \widehat\Phi:
        M\times P
        \longrightarrow
        K(P\times N)
    \]
    satisfying
    \begin{equation}
        \widehat\Phi(e_M,x)
        =
        \eta^K(x,e_N),
        \label{eq:one-object-effectful-unit}
    \end{equation}
    and
    \begin{equation}
    \begin{aligned}
        \widehat\Phi(mn,x)
        =
        \operatorname{bind}_{K}
        \Bigl(
            \widehat\Phi(m,x),
            (x_1,u)
            \longmapsto
            \operatorname{bind}_{K}
            \bigl(
                \widehat\Phi(n,x_1),
                (x_2,v)
                \longmapsto
                \eta^K(x_2,uv)
            \bigr)
        \Bigr).
    \end{aligned}
    \label{eq:one-object-effectful-composition}
    \end{equation}
\end{proposition}
\begin{proof}
    The carrier span reduces to the set \(P\), while
    \[
        T_N(P)=P\times N
    \]
    and
    \[
        R_M(P)=M\times P.
    \]
    Proposition~\ref{prop:effectful-mealy-laws} therefore applies.

    Under the one-object convention, an outcome of the first execution is
    displayed as
    \[
        \begin{tikzcd}[column sep=huge]
                x
                    \arrow[r, "{m/u}"]
            &
                x_1,
        \end{tikzcd}
    \]
    and an outcome of the second execution is
    \[
        \begin{tikzcd}[column sep=huge]
                x_1
                    \arrow[r, "{n/v}"]
            &
                x_2.
        \end{tikzcd}
    \]
    The two outputs compose to \(uv\), since categorical composition in the
    one-object output category satisfies
    \[
        v\circ u=uv.
    \]
    Hence the general fibrewise law becomes
    \eqref{eq:one-object-effectful-composition}.
\end{proof}

When \(K=\mathrm{Id}\), write
\[
    \widehat\Phi(m,x)
    =
    (x\cdot m,\omega(m,x)).
\]
Then
\eqref{eq:one-object-effectful-composition} reduces to
\[
    x\cdot(mn)
    =
    (x\cdot m)\cdot n
\]
and
\[
    \omega(mn,x)
    =
    \omega(m,x)\omega(n,x\cdot m).
\]

\subsection{Effectful Mealy simulations}
\label{subsec:effectful-mealy-simulations}

Let
\[
    (P,\widehat\Phi),
    (Q,\widehat\Psi):
    \mathbb A\rightsquigarrow_{\mathsf K}\mathbb B
\]
be effectful Mealy morphisms. They are algebras for
\[
    \overline R_A^{\mathsf K}
\]
in
\[
    \operatorname{Kl}(H_B^{\mathsf K}).
\]

As in the deterministic theory, the displayed simulation direction is chosen
opposite to the underlying algebra-morphism direction.

\begin{definition}[Effectful Mealy simulation]
\label{def:effectful-mealy-simulation}
    An \textbf{effectful Mealy simulation}
    \[
        (P,\widehat\Phi)
        \Rightarrow
        (Q,\widehat\Psi)
    \]
    is a morphism
    \[
        (Q,\widehat\Psi)
        \longrightarrow
        (P,\widehat\Phi)
    \]
    of \(\overline R_A^{\mathsf K}\)-algebras in
    \[
        \operatorname{Kl}(H_B^{\mathsf K}).
    \]

    Equivalently, it consists of a map
    \[
        \widehat\Theta:
        Q
        \longrightarrow
        H_B^{\mathsf K}(P)
        =
        \mathsf K T_B(P)
    \]
    satisfying
    \begin{equation}
    \begin{aligned}
        &
        \mu^H_P
        \circ
        H_B^{\mathsf K}(\widehat\Theta)
        \circ
        \widehat\Psi
        \\
        &\qquad=
        \mu^H_P
        \circ
        H_B^{\mathsf K}(\widehat\Phi)
        \circ
        \chi_P
        \circ
        R_A(\widehat\Theta).
    \end{aligned}
    \label{eq:effectful-simulation-abstract}
    \end{equation}
\end{definition}

Thus a state of \(Q\) may be compared with an effectful value whose pure
outcomes consist of a state of \(P\) and an output-comparison arrow:
\[
    \begin{tikzcd}[column sep=huge]
            q_P(x)
                \arrow[r, "\alpha"]
        &
            q_Q(y).
    \end{tikzcd}
\]

A useful subclass retains deterministic comparison data.

\begin{definition}[Pure effectful Mealy simulation]
\label{def:pure-effectful-mealy-simulation}
    A \textbf{pure effectful Mealy simulation}
    \[
        (P,\widehat\Phi)
        \Rightarrow
        (Q,\widehat\Psi)
    \]
    consists of a map of spans
    \[
        \Theta:
        Q
        \Rightarrow
        T_B(P)
    \]
    such that
    \[
        \widehat\Theta
        :=
        j_P\circ\Theta
        =
        \eta^{\mathsf K}_{T_B(P)}
        \circ
        \Theta
    \]
    is an effectful Mealy simulation.

    Writing
    \[
        \Theta(y)
        =
        (\theta y,\alpha_y),
    \]
    its component data satisfy
    \[
        p_P\theta=p_Q,
    \]
    \[
        s_B\alpha=q_P\theta,
    \]
    and
    \[
        t_B\alpha=q_Q.
    \]
\end{definition}

The comparison data are displayed as
\[
    \begin{tikzcd}[column sep=large, row sep=large]
            &
                Q
                    \arrow[d, "\theta" description]
                    \arrow[ddl, "p_Q"', bend right=15]
                    \arrow[ddr, "q_Q", bend left=15]
                    \arrow[dr, "\alpha"]
            &
        \\
            &
                P
                    \arrow[dl, "p_P"]
                    \arrow[dr, "q_P"']
            &
                B_1
                    \arrow[d, "s_B"]
                    \arrow[ddr, "t_B", bend left=18]
        \\
                A_0
            &
            &
                B_0
        \\
            &
            &
            &
                B_0.
    \end{tikzcd}
\]

\subsection{Fibrewise pure simulation laws}
\label{subsec:fibrewise-pure-effectful-simulation-laws}

Assume again that
\[
    \mathcal E=\mathbf{Set}
\]
and that the effect datum is induced fibrewise by a Set-monad \(K\).

\begin{proposition}[The fibrewise pure effectful simulation law]
\label{prop:pure-effectful-simulation-law}
    A pair
    \[
        (\theta,\alpha)
    \]
    as in Definition~\ref{def:pure-effectful-mealy-simulation} is a pure
    effectful Mealy simulation if and only if, for every admissible
    input-state pair
    \[
        (a,y),
    \]
    one has
    \begin{equation}
    \begin{aligned}
        &
        \operatorname{bind}_{K}
        \Bigl(
            \widehat\Phi(a,\theta y),
            (x,\beta)
            \longmapsto
            \eta^K
            (x,\alpha_y\circ\beta)
        \Bigr)
        \\
        &\qquad=
        \operatorname{bind}_{K}
        \Bigl(
            \widehat\Psi(a,y),
            (y',\gamma)
            \longmapsto
            \eta^K
            \bigl(
                \theta y',
                \gamma\circ\alpha_{y'}
            \bigr)
        \Bigr).
    \end{aligned}
    \label{eq:pure-effectful-simulation-component}
    \end{equation}
\end{proposition}
\begin{proof}
    \leavevmode
    \begin{enumerate}
        \item \textbf{Compare first and then execute.}
        The state \(y\) is first compared with \(\theta y\) by
        \[
            \alpha_y:
            q_P(\theta y)\to q_Q(y).
        \]
        Executing \(P\) along \(a\) at \(\theta y\) produces effectfully a
        pair
        \[
            (x,\beta)
        \]
        with
        \[
            \beta:
            q_P(x)\to q_P(\theta y).
        \]
        The resulting comparison arrow is
        \[
            \begin{tikzcd}[column sep=huge]
                    q_P(x)
                        \arrow[r, "\beta"]
                &
                    q_P(\theta y)
                        \arrow[r, "\alpha_y"]
                &
                    q_Q(y).
            \end{tikzcd}
        \]
        Hence the output component is
        \[
            \alpha_y\circ\beta.
        \]

        \item \textbf{Execute first and then compare.}
        Executing \(Q\) along \(a\) at \(y\) produces effectfully
        \[
            (y',\gamma)
        \]
        with
        \[
            \gamma:
            q_Q(y')\to q_Q(y).
        \]
        The updated state is compared with \(\theta y'\) by
        \[
            \alpha_{y'}:
            q_P(\theta y')\to q_Q(y').
        \]
        The resulting output arrow is
        \[
            \begin{tikzcd}[column sep=huge]
                    q_P(\theta y')
                        \arrow[r, "\alpha_{y'}"]
                &
                    q_Q(y')
                        \arrow[r, "\gamma"]
                &
                    q_Q(y).
            \end{tikzcd}
        \]
        Hence the output component is
        \[
            \gamma\circ\alpha_{y'}.
        \]

        \item \textbf{Compare the two effectful outcomes.}
        The algebra-morphism equation
        \eqref{eq:effectful-simulation-abstract} says precisely that these two
        effectful joint state-output values agree. This is
        \eqref{eq:pure-effectful-simulation-component}.
    \end{enumerate}
\end{proof}

\begin{remark}[Recovery of deterministic simulations]
\label{rem:effectful-simulations-recover-deterministic}
    If
    \[
        K=\mathrm{Id},
    \]
    then
    \eqref{eq:pure-effectful-simulation-component} is equality of the pairs
    \[
        \bigl(
            \delta_P(a,\theta y),
            \alpha_y\circ\omega_P(a,\theta y)
        \bigr)
    \]
    and
    \[
        \bigl(
            \theta(\delta_Q(a,y)),
            \omega_Q(a,y)\circ
            \alpha_{\delta_Q(a,y)}
        \bigr).
    \]
    Equality of the state and output components recovers
    \[
        \theta(\delta_Q(a,y))
        =
        \delta_P(a,\theta y)
    \]
    and
    \[
        \alpha_y\circ\omega_P(a,\theta y)
        =
        \omega_Q(a,y)\circ
        \alpha_{\delta_Q(a,y)}.
    \]
\end{remark}

\begin{proposition}[Vertical composition of effectful simulations]
\label{prop:vertical-composition-effectful-simulations}
    Effectful Mealy simulations compose vertically.

    Pure effectful simulations are closed under vertical composition. If
    \[
        (\theta,\alpha):
        P\Rightarrow Q
    \]
    and
    \[
        (\kappa,\beta):
        Q\Rightarrow R
    \]
    are pure effectful simulations, then their composite is pure, with state
    map
    \[
        \theta\kappa:R\to P
    \]
    and comparison component
    \[
        \gamma_z
        =
        \beta_z\circ\alpha_{\kappa z}.
    \]
\end{proposition}
\begin{proof}
    General effectful simulations are morphisms in the category of
    \(\overline R_A^{\mathsf K}\)-algebras in
    \(\operatorname{Kl}(H_B^{\mathsf K})\), so their vertical composition is
    inherited from that category.

    The monad morphism
    \[
        j:T_B\to H_B^{\mathsf K}
    \]
    induces a functor of Kleisli categories. It therefore preserves Kleisli
    composition. Hence the composite of two pure simulations is the image
    under \(j\) of their composite in \(\operatorname{Kl}(T_B)\).

    Componentwise, the comparison arrows form the path
    \[
        \begin{tikzcd}[column sep=huge]
                q_P(\theta\kappa z)
                    \arrow[r, "\alpha_{\kappa z}"]
            &
                q_Q(\kappa z)
                    \arrow[r, "\beta_z"]
            &
                q_R(z),
        \end{tikzcd}
    \]
    whose composite is
    \[
        \beta_z\circ\alpha_{\kappa z}.
    \]
\end{proof}

\subsection{Effectful pipeline composition}
\label{subsec:effectful-pipeline-composition}

Effectful pipeline composition requires more than a monad on
\(\mathbf{Set}\). The effects generated by different stages of a pipeline must
be interchangeable when they are conditionally independent. We therefore
assume in this subsection that
\[
    K=(K,\eta^K,\mu^K)
\]
is a \textbf{commutative monad} on \(\mathbf{Set}\), and use its canonical
fibrewise effect doctrine.

The finite-powerset, full-powerset, finite-distribution,
finite-subdistribution, and lift monads are commutative. The weighted multiset
monad associated to a semiring is commutative when the semiring is
commutative.

Let
\[
    (P,\widehat\Phi_P):
    \mathbb A\rightsquigarrow_K\mathbb B
\]
and
\[
    (Q,\widehat\Phi_Q):
    \mathbb B\rightsquigarrow_K\mathbb C
\]
be effectful Mealy morphisms. Their composite carrier is the ordinary
pullback carrier
\[
    P\times_{q_P,B_0,p_Q}Q.
\]
A composite state is a pair
\[
    (x,y)
\]
satisfying
\[
    q_P(x)=p_Q(y).
\]

\begin{proposition}[Effectful pipeline composition in
\texorpdfstring{\(\mathbf{Set}\)}{Set}]
\label{prop:effectful-pipeline-composition-set}
    Let \(K\) be a commutative monad on \(\mathbf{Set}\). The formula
    \begin{equation}
    \begin{aligned}
        &
        \widehat\Phi_{Q\circ P}(a,(x,y))
        \\
        &\quad=
        \operatorname{bind}_{K}
        \Bigl(
            \widehat\Phi_P(a,x),
            (x',\beta)
            \longmapsto
            \operatorname{bind}_{K}
            \bigl(
                \widehat\Phi_Q(\beta,y),
                (y',\gamma)
                \longmapsto
                \eta^K
                \bigl(
                    (x',y'),
                    \gamma
                \bigr)
            \bigr)
        \Bigr)
    \end{aligned}
    \label{eq:effectful-pipeline-composition}
    \end{equation}
    defines an effectful Mealy morphism
    \[
        Q\circ P:
        \mathbb A\rightsquigarrow_K\mathbb C.
    \]
\end{proposition}
\begin{proof}
    \leavevmode
    \begin{enumerate}
        \item \textbf{Check the intermediate typing.}
        An outcome of the first machine has the form
        \[
            (x',\beta)
        \]
        with
        \[
            \beta:
            q_P(x')\to q_P(x).
        \]
        Since
        \[
            q_P(x)=p_Q(y),
        \]
        the arrow \(\beta\) is admissible input for \(Q\) at \(y\):
        \[
            \begin{tikzcd}[column sep=huge]
                    q_P(x')
                        \arrow[r, "\beta"]
                &
                    q_P(x)=p_Q(y).
            \end{tikzcd}
        \]

        \item \textbf{Check the updated carrier.}
        An outcome of executing \(Q\) along \(\beta\) at \(y\) has the form
        \[
            (y',\gamma)
        \]
        with
        \[
            p_Q(y')
            =
            s_B(\beta)
            =
            q_P(x').
        \]
        Hence
        \[
            (x',y')
            \in
            P\times_{B_0}Q.
        \]

        \item \textbf{Check the output typing.}
        The output arrow is
        \[
            \gamma:
            q_Q(y')\to q_Q(y),
        \]
        which has exactly the required type for the composite carrier.

        The execution is represented by
        \[
            \begin{tikzcd}[column sep=huge, row sep=large]
                    x
                        \arrow[r, "{a/\beta}"]
                &
                    x'
                \\
                    y
                        \arrow[r, "{\beta/\gamma}"']
                &
                    y'.
            \end{tikzcd}
        \]

        \item \textbf{Verify the unit law.}
        On the identity input \(e_A(p_Px)\), the first machine produces the
        pure outcome
        \[
            (x,e_B(q_Px)).
        \]
        The second machine then produces the pure outcome
        \[
            (y,e_C(q_Qy)).
        \]
        The monad unit equations reduce
        \eqref{eq:effectful-pipeline-composition} to
        \[
            \eta^K
            \bigl(
                (x,y),
                e_C(q_Qy)
            \bigr).
        \]

        \item \textbf{Expand sequential execution.}
        Let
        \[
            a:u\to v,
            \qquad
            b:v\to w,
            \qquad
            p_P(x)=w.
        \]
        Sequential execution of \(b\) and then \(a\) in the pipeline has the
        form
        \[
        \begin{aligned}
            &
            \operatorname{bind}_{K}
            \Bigl(
                \widehat\Phi_P(b,x),
                (x_b,\beta_b)
                \longmapsto
                \operatorname{bind}_{K}
                \Bigl(
                    \widehat\Phi_Q(\beta_b,y),
                    (y_b,\gamma_b)
                    \longmapsto
            \\
            &\hspace{6em}
                    \operatorname{bind}_{K}
                    \Bigl(
                        \widehat\Phi_P(a,x_b),
                        (x_a,\beta_a)
                        \longmapsto
                        \operatorname{bind}_{K}
                        \bigl(
                            \widehat\Phi_Q(\beta_a,y_b),
                            (y_a,\gamma_a)
            \\
            &\hspace{14em}
                            \longmapsto
                            \eta^K
                            \bigl(
                                (x_a,y_a),
                                \gamma_b\circ\gamma_a
                            \bigr)
                        \bigr)
                    \Bigr)
                \Bigr)
            \Bigr).
        \end{aligned}
        \label{eq:effectful-pipeline-sequential-expanded}
        \]

        \item \textbf{Expand direct execution.}
        The multiplication law for \(P\) expands
        \[
            \widehat\Phi_P(b\circ a,x)
        \]
        into the sequence
        \[
            \widehat\Phi_P(b,x),
            \qquad
            \widehat\Phi_P(a,x_b),
        \]
        with emitted composite arrow
        \[
            \beta_b\circ\beta_a.
        \]
        The multiplication law for \(Q\) then expands
        \[
            \widehat\Phi_Q(\beta_b\circ\beta_a,y)
        \]
        into the sequence
        \[
            \widehat\Phi_Q(\beta_b,y),
            \qquad
            \widehat\Phi_Q(\beta_a,y_b),
        \]
        with output
        \[
            \gamma_b\circ\gamma_a.
        \]

        Before using commutativity, the direct expansion orders its middle
        computations as
        \[
            \widehat\Phi_P(a,x_b)
            \quad\text{then}\quad
            \widehat\Phi_Q(\beta_b,y),
        \]
        whereas sequential pipeline execution orders them as
        \[
            \widehat\Phi_Q(\beta_b,y)
            \quad\text{then}\quad
            \widehat\Phi_P(a,x_b).
        \]

        \item \textbf{Use commutativity of the effect.}
        After the first outcome
        \[
            (x_b,\beta_b)
        \]
        has been fixed, the computations
        \[
            \widehat\Phi_P(a,x_b)
        \]
        and
        \[
            \widehat\Phi_Q(\beta_b,y)
        \]
        are conditionally independent. Commutativity of \(K\) interchanges
        these two computations without changing their joint effectful value.

        After this interchange, associativity of Kleisli extension identifies
        the direct expansion with
        \eqref{eq:effectful-pipeline-sequential-expanded}. Associativity of
        composition in \(\mathbb C\) identifies the emitted output arrows.
        Hence the multiplication law holds.
    \end{enumerate}
\end{proof}

\begin{remark}[Why commutativity is required]
\label{rem:effectful-pipeline-commutativity-required}
    For a general noncommutative Set-monad, direct execution of a composite
    input and sequential execution through a pipeline can order independent
    effects differently.

    After the first \(P\)-outcome has been selected, the direct expansion
    orders the middle computations as
    \[
        P(a)
        \quad\text{then}\quad
        Q(\beta_b),
    \]
    while sequential pipeline execution orders them as
    \[
        Q(\beta_b)
        \quad\text{then}\quad
        P(a).
    \]
    A commutative monad supplies the interchange equating these two orders.
    Without such interchange, formula
    \eqref{eq:effectful-pipeline-composition} need not satisfy the effectful
    Mealy multiplication law.
\end{remark}

\begin{remark}[Associativity and simulations]
\label{rem:associativity-effectful-pipelines}
    For fibrewise effects induced by a commutative Set-monad, effectful
    pipeline composition is associative up to the canonical pullback
    associator
    \[
        (P\times_{B_0}Q)\times_{C_0}R
        \cong
        P\times_{B_0}(Q\times_{C_0}R).
    \]
    The proof uses associativity and commutativity of Kleisli extension,
    together with the canonical associativity isomorphism of pullbacks.

    The identity effectful Mealy morphism is the pure effectful morphism
    obtained from the deterministic identity Mealy morphism.

    General effectful simulations compose vertically because they are
    algebra morphisms. Constructing horizontal composition of general
    effectful simulations requires an additional calculation compatible with
    effectful pipeline composition and is not asserted here.

    Pure simulations are closed under vertical composition, but they are not
    generally closed under horizontal composition. Horizontally composing two
    pure comparisons may require executing an effectful machine on an
    output-comparison arrow, and the resulting comparison is then effectful
    rather than pure.
\end{remark}

\subsection{Nondeterministic Mealy morphisms}
\label{subsec:nondeterministic-mealy-morphisms}

Take
\[
    K=\mathcal P_{\mathrm f},
\]
the finite-powerset monad. An effectful transition is then a finite set of
possible next-state/output pairs:
\[
    \widehat\Phi(a,x)
    \subseteq
    \operatorname{Out}_P(s_A(a),q(x)).
\]

\begin{example}[Nondeterministic Mealy morphisms]
\label{example:nondeterministic-mealy-morphisms}
    A finite-nondeterministic Mealy morphism
    \[
        \mathbb A
        \rightsquigarrow_{\mathcal P_{\mathrm f}}
        \mathbb B
    \]
    consists of a carrier span
    \[
        A_0\xleftarrow{p}P\xrightarrow{q}B_0
    \]
    and a finite set
    \[
        \widehat\Phi(a,x)
    \]
    of typed outcomes for every admissible input-state pair.

    A branch has the form
    \[
        \begin{tikzcd}[column sep=huge, row sep=large]
                u
                    \arrow[r, "a"]
            &
                v=p(x)
            \\
                q(x')
                    \arrow[r, "\beta"']
            &
                q(x).
        \end{tikzcd}
    \]

    The unit law is
    \[
        \widehat\Phi(e_A(px),x)
        =
        \bigl\{
            (x,e_B(qx))
        \bigr\}.
    \]

    The multiplication law is
    \[
    \begin{aligned}
        \widehat\Phi(b\circ a,x)
        =
        \bigl\{
            (x_a,\beta_b\circ\beta_a)
            \ \big|\
            &(x_b,\beta_b)\in\widehat\Phi(b,x),
            \\
            &(x_a,\beta_a)\in\widehat\Phi(a,x_b)
        \bigr\}.
    \end{aligned}
    \]
    Thus the outcomes of a composite input are exactly the outcomes obtained
    by choosing one first-step branch and then one second-step branch.

    In the one-object case, this is a map
    \[
        M\times P
        \longrightarrow
        \mathcal P_{\mathrm f}(P\times N)
    \]
    satisfying
    \[
        \widehat\Phi(1,x)
        =
        \{(x,1)\}
    \]
    and
    \[
    \begin{aligned}
        \widehat\Phi(mn,x)
        =
        \bigl\{
            (x_2,uv)
            \ \big|\
            &(x_1,u)\in\widehat\Phi(m,x),
            \\
            &(x_2,v)\in\widehat\Phi(n,x_1)
        \bigr\}.
    \end{aligned}
    \]
\end{example}

The full powerset monad allows arbitrary branching. The nonempty powerset
monad excludes deadlock, while the ordinary powerset monad permits the empty
set of outcomes.

For input and output alphabets \(\Sigma\) and \(\Gamma\), a nondeterministic
classical Mealy machine is obtained by taking
\[
    M=\Sigma^\ast,
    \qquad
    N=\Gamma^\ast,
\]
and requiring generator-level outcomes to emit single letters when the
classical letter-output convention is desired.

\begin{remark}[Strictness of nondeterministic simulations]
\label{rem:strictness-nondeterministic-simulations}
    For the finite-powerset effect, the effectful simulation equation is an
    equality of finite sets of compared outcomes. Hence the simulations
    defined above are strict Kleisli algebra morphisms.

    A one-sided relational simulation would instead use a lax
    algebra-morphism condition, expressed by inclusion in the
    order-enriched Kleisli category of the powerset monad. Such lax
    simulations are a distinct enrichment of the present theory.
\end{remark}

\subsection{Probabilistic Mealy morphisms}
\label{subsec:probabilistic-mealy-morphisms}

Let
\[
    \mathcal D_{\mathrm f}(X)
\]
denote the set of finitely supported probability distributions on \(X\). Its
unit is the Dirac distribution, and its multiplication integrates a
distribution of distributions.

\begin{example}[Probabilistic Mealy morphisms]
\label{example:probabilistic-mealy-morphisms}
    Take
    \[
        K=\mathcal D_{\mathrm f}.
    \]
    A probabilistic Mealy morphism
    \[
        \mathbb A
        \rightsquigarrow_{\mathcal D_{\mathrm f}}
        \mathbb B
    \]
    assigns to every admissible pair \((a,x)\) a finitely supported
    probability distribution
    \[
        \widehat\Phi(a,x)
    \]
    on the typed next-state/output pairs
    \[
        (x',\beta),
        \qquad
        \beta:q(x')\to q(x).
    \]

    The identity input has the Dirac distribution
    \[
        \widehat\Phi(e_A(px),x)
        =
        \delta_{(x,e_B(qx))}.
    \]

    For composable arrows
    \[
        a:u\to v,
        \qquad
        b:v\to w,
    \]
    the probability of obtaining the outcome
    \[
        (x_a,\gamma)
    \]
    after processing \(b\circ a\) is
    \[
    \begin{aligned}
        &
        \widehat\Phi(b\circ a,x)(x_a,\gamma)
        \\
        &\quad=
        \sum_{\substack{
            x_b,\beta_b,\beta_a\\
            \beta_b\circ\beta_a=\gamma
        }}
        \widehat\Phi(b,x)(x_b,\beta_b)\,
        \widehat\Phi(a,x_b)(x_a,\beta_a).
    \end{aligned}
    \]
    Only supported pairs contribute, so the sum is finite even when the
    output category admits infinitely many abstract factorizations of
    \(\gamma\).

    The output factorization is
    \[
        \begin{tikzcd}[column sep=huge]
                q(x_a)
                    \arrow[r, "\beta_a"]
            &
                q(x_b)
                    \arrow[r, "\beta_b"]
            &
                q(x).
        \end{tikzcd}
    \]
    Thus the effectful multiplication law is a typed
    Chapman--Kolmogorov convolution: probabilities multiply along sequential
    executions and are summed over hidden intermediate states and output-arrow
    factorizations.

    In the one-object case, a probabilistic Mealy morphism is a finite-support
    Markov kernel
    \[
        M\times P
        \longrightarrow
        \mathcal D_{\mathrm f}(P\times N)
    \]
    satisfying
    \[
        \widehat\Phi(1,x)
        =
        \delta_{(x,1)}
    \]
    and
    \[
    \begin{aligned}
        \widehat\Phi(mn,x)(x_2,w)
        =
        \sum_{\substack{x_1,u,v\\uv=w}}
        \widehat\Phi(m,x)(x_1,u)\,
        \widehat\Phi(n,x_1)(x_2,v).
    \end{aligned}
    \]
\end{example}

The distribution is placed on \(P\times N\), rather than separately on \(P\)
and \(N\). Hence the next state and output may be probabilistically correlated.
A model in which state and output are conditionally independent is the special
case in which the joint distribution factors.

Replacing \(\mathcal D_{\mathrm f}\) by the finite-subdistribution monad
allows loss of probability mass and therefore models failure, rejection, or
nontermination in addition to probabilistic choice.

\subsection{Partial and weighted Mealy morphisms}
\label{subsec:partial-weighted-mealy-morphisms}

\begin{example}[Partial Mealy morphisms]
\label{example:partial-mealy-morphisms}
    Take the lift monad
    \[
        K(X)=X+1.
    \]
    An effectful transition
    \[
        \widehat\Phi(a,x)
    \]
    is either a successful typed next-state/output pair
    \[
        \operatorname{inl}(x',\beta)
    \]
    or the distinguished failure element
    \[
        \operatorname{inr}(\ast).
    \]

    The unit law requires identity inputs to succeed:
    \[
        \widehat\Phi(e_A(px),x)
        =
        \operatorname{inl}
        \bigl(x,e_B(qx)\bigr).
    \]

    Sequential execution fails if either stage fails. If both stages succeed
    with outcomes
    \[
        (x_b,\beta_b)
        \qquad\text{and}\qquad
        (x_a,\beta_a),
    \]
    then the composite succeeds with
    \[
        (x_a,\beta_b\circ\beta_a).
    \]

    The successful output path is
    \[
        \begin{tikzcd}[column sep=huge]
                q(x_a)
                    \arrow[r, "\beta_a"]
            &
                q(x_b)
                    \arrow[r, "\beta_b"]
            &
                q(x).
        \end{tikzcd}
    \]

    In the one-object case this is a partial function
    \[
        M\times P
        \rightharpoonup
        P\times N
    \]
    compatible with the monoid multiplication.
\end{example}

\begin{example}[Semiring-weighted Mealy morphisms]
\label{example:semiring-weighted-mealy-morphisms}
    Let \(R\) be a semiring, and let
    \[
        \mathsf W_R(X)
    \]
    be the set of finitely supported functions
    \[
        X\to R.
    \]
    The unit assigns weight \(1\) to one point and \(0\) elsewhere, while
    multiplication sums products of weights over intermediate points.

    An \(\mathsf W_R\)-effectful Mealy morphism assigns an \(R\)-weight to
    every possible typed next-state/output pair. Sequential execution
    multiplies the weights of consecutive outcomes and sums over intermediate
    states and output-arrow factorizations.

    If
    \[
        w_{a,x}(x',\beta)\in R
    \]
    denotes the weight assigned to an outcome, then
    \[
    \begin{aligned}
        w_{b\circ a,x}(x_a,\gamma)
        =
        \sum_{\substack{
            x_b,\beta_b,\beta_a\\
            \beta_b\circ\beta_a=\gamma
        }}
        w_{b,x}(x_b,\beta_b)\,
        w_{a,x_b}(x_a,\beta_a).
    \end{aligned}
    \]

    The Boolean semiring recovers finite nondeterminism. The natural-number
    semiring counts executions with multiplicity. Other semirings yield
    weighted transducers internal to the typed input-output calculus.

    When \(R\) is commutative, the monad \(\mathsf W_R\) is commutative and
    therefore supports the pipeline composition of
    Proposition~\ref{prop:effectful-pipeline-composition-set}.
\end{example}

\subsection{Relation to the deterministic theory}
\label{subsec:effectful-relation-deterministic-theory}

The effectful construction extends the deterministic theory in three
compatible ways.

First, taking
\[
    \mathsf K=\mathrm{Id}
\]
gives
\[
    H_B^{\mathsf K}=T_B,
\]
and Definition~\ref{def:effectful-mealy-morphism} reduces to
Definition~\ref{def:mealy-morphism}.

Second, for an arbitrary Mealy-compatible effect datum, every deterministic
Mealy morphism embeds as a pure effectful Mealy morphism by
Proposition~\ref{prop:pure-effectful-mealy-morphisms}. The embedding is induced
by the monad morphism
\[
    j:
    T_B
    \longrightarrow
    \mathsf K T_B.
\]

Third, pure effectful simulations retain the deterministic state-comparison
and output-comparison data of ordinary Mealy simulations. Their compatibility
equation is no longer generally a pointwise equality of selected transitions;
it is equality of the two induced effectful joint state-output behaviours.

The principal fibrewise Set effects may be summarized as follows:
\[
\begin{array}{c|l}
\text{effect monad} & \text{resulting Mealy behaviour} \\
\hline
\mathrm{Id}
&
\text{deterministic state transition and output}
\\
\mathcal P_{\mathrm f}
&
\text{finite nondeterministic branching}
\\
\mathcal P
&
\text{arbitrary nondeterministic branching}
\\
\mathcal D_{\mathrm f}
&
\text{finitely supported probabilistic branching}
\\
\text{finite subdistribution}
&
\text{probability with failure or nontermination}
\\
(-)+1
&
\text{partial transition and output}
\\
\mathsf W_R
&
\text{semiring-weighted transition and output.}
\end{array}
\]

In each case, the effect is applied to the joint next-state/output object
\[
    B\circ P.
\]
The structural sequence is
\[
    \begin{tikzcd}[column sep=huge]
            P\circ A
                \arrow[r, "\widehat\Phi"]
        &
            \mathsf K(B\circ P),
    \end{tikzcd}
\]
rather than a pair of separate maps selecting an effectful state and an
effectful output.

The input-processing monad acts through the distributive law
\[
    \begin{tikzcd}[column sep=huge]
            R_A\mathsf K T_B
                \arrow[r, "\rho_A T_B"]
        &
            \mathsf K R_AT_B
                \arrow[r, "\mathsf K\lambda"]
        &
            \mathsf K T_BR_A.
    \end{tikzcd}
\]
Consequently, identity and sequential execution remain algebra laws for a
lifted input-processing monad. This preserves the structural organization of
the deterministic theory while allowing the transition-output operation to
carry computational effects.

For commutative fibrewise Set effects, the pipeline formula further composes
effectful Mealy morphisms:
\[
    \begin{tikzcd}[column sep=large, row sep=large]
            x
                \arrow[r, "{a/\beta}"]
        &
            x'
        \\
            y
                \arrow[r, "{\beta/\gamma}"']
        &
            y'.
    \end{tikzcd}
\]
The emitted arrow of the first machine becomes the effectful input of the
second, while the two computational effects are combined by Kleisli extension
and commutative interchange.